\section{Antecedentes Empíricos}

La literatura sobre las dinámicas distributivas, el costo de vida y los mercados laborales en América Latina evidencia un cuerpo amplio de investigaciones aplicadas. Para abordar la relación entre los ingresos de los hogares y el costo de vida en Chile, Brasil y Argentina, la evidencia empírica seleccionada se sistematiza en cinco ejes analíticos fundamentales: i) estructura laboral, política salarial e institucionalidad del trabajo; ii) dinámica inflacionaria, formación de precios y costo de vida; iii) protección social, transferencias monetarias y mitigación de la pobreza; iv) concentración de ingresos, riqueza y brechas interseccionales (raza y género); y v) movilidad social, estratificación y percepciones de la desigualdad. La definición de estos ejes se fundamenta en la revisión de la literatura especializada y su proceso de sistematización detallado a continuación.

\subsection{Estrategia de Búsqueda, Limitaciones y Fuentes Institucionales}

Para fundamentar el estado del arte y garantizar un marco de antecedentes transparente, exhaustivo y replicable, se diseñó una estrategia de búsqueda bibliográfica híbrida basada en los lineamientos del protocolo PRISMA 2020\footnote[2]{PRISMA significa \textit{Preferred Reporting Items for Systematic Reviews and Meta-Analyses} (Elementos de información preferidos para revisiones sistemáticas y metaanálisis). La versión 2020 reemplaza a la antigua guía de 2009 para adaptarse a los nuevos avances metodológicos en la investigación. https://www.sciencedirect.com/science/article/pii/S0300893221002748} (Figura~\ref{fig:prisma-flowchart}) se profundizan en el \textbf{Anexo 1}. Esta estrategia articuló la captura sistemática de literatura científica de alto impacto en la colección principal de \textit{Web of Science} (WoS) con búsquedas complementarias en repositorios regionales de acceso abierto (\textit{SciELO}, \textit{Redalyc}, \textit{Dialnet} y \textit{Google Académico}) en español y portugués, sumada a la incorporación estructurada de literatura gris institucional.\footnote{Se entiende por literatura gris aquel conjunto de documentos y publicaciones técnicas de carácter público u oficial —como informes de organismos multilaterales, boletines estadísticos, documentos de trabajo y memorias metodológicas— que no son distribuidos a través de los canales de edición comercial ni de revistas científicas sujetas a revisión por pares tradicional.}

Si bien este abordaje permite capturar la producción académica más relevante sobre la intersección entre ingresos, costo de vida e inflación en el Cono Sur, la literatura científica puramente indexada en plataformas internacionales presenta tres limitaciones metodológicas clave:

\begin{itemize}
    \item \textbf{Sesgo de indexación bibliométrica e idioma}: Las bases de datos internacionales \textit{mainstream} priorizan publicaciones en idioma inglés, corriendo el riesgo de omitir valiosos aportes empíricos y discusiones teóricas editadas en revistas latinoamericanas de ámbito regional.
    \item \textbf{Desfase temporal de publicación (\textit{Publication Lag})}: El proceso de arbitraje por pares impone una latencia editorial (frecuentemente superior a los 12--18 meses) entre el levantamiento de datos y la publicación final, lo que resta agilidad para examinar los efectos inmediatos de shocks inflacionarios y coyunturas macroeconómicas recientes (2021--2025).
    \item \textbf{Nivel de agregación de los análisis}: La producción \textit{peer-reviewed} tiende a enfocarse en modelos teóricos o variables macroeconómicas agregadas, ofreciendo escasa documentación técnica sobre el procesamiento operativo y la armonización directa de microdatos de encuestas de hogares.
\end{itemize}

Por estos motivos, resultó indispensable adoptar un enfoque híbrido que complemente los artículos científicos con la literatura gris institucional producida por organismos multilaterales (\textit{CEPAL}) e institutos de estadística oficiales (\textit{INDEC} en Argentina, \textit{IBGE} en Brasil e \textit{INE/MIDES} en Chile). La integración sistemática de esta literatura gris se fundamenta en tres pilares:

\begin{enumerate}
    \item \textbf{Actualización y oportunidad de datos}: Acceso a mediciones mensuales y trimestrales del Índice de Precios al Consumidor (IPC), la Canasta Básica Alimentaria (CBA) y la Canasta Básica Total (CBT), evitando los rezagos del ciclo editorial académico.
    \item \textbf{Precisión metodológica oficial}: Disponibilidad de la documentación técnica sobre el diseño muestral, criterios de imputación y metodologías de deflactación de las encuestas oficiales de hogares (\textit{CASEN}, \textit{PNAD Contínua} y \textit{EPH}).
    \item \textbf{Estándares de comparabilidad regional}: Criterios de armonización monetaria y Paridad de Poder Adquisitivo (PPA) promovidos por la CEPAL en publicaciones como el \textit{Panorama Social de América Latina y el Caribe}.
\end{enumerate}

El procesamiento de metadatos se ejecutó de forma algorítmica en lenguaje \textit{R} mediante el paquete \texttt{bibliometrix}. La búsqueda inicial en \textit{Web of Science} arrojó un total de 3.982 registros brutos. Tras remover 18 elementos duplicados, se evaluaron automatizadamente 3.964 artículos por título y resumen. La aplicación de filtros de expresiones regulares (\texttt{grep}) excluyó 218 registros no pertinentes, consolidando un corpus cribado de 3.746 publicaciones. Este corpus fue importado al gestor bibliográfico \textit{Zotero}, donde se llevó a cabo una etapa de curaduría cualitativa y selección manual estructurada mediante etiquetado temático (\textit{tags}) por ejes analíticos y revisión a texto completo en una subcarpeta final de inclusión. Durante este análisis de elegibilidad se descartaron 3.241 registros, consolidando una muestra definitiva de **104 estudios incluidos** para el documento.

El desglose detallado de las cadenas de consulta bibliográfica, la matriz completa de criterios de inclusión y exclusión, la parametrización del filtrado algorítmico y el diagrama de flujo metodológico estructurado bajo el protocolo PRISMA 2020 se detallan de manera integral en el \textbf{Anexo~\ref{sec:procesamiento_algoritmico_filtrado}} (véase la Figura~\ref{fig:prisma-flowchart}).

\subsection{i. Mercado de trabajo, instituciones salariales y distribución del ingreso laboral}

La dinámica de las remuneraciones en la región se encuentra fuertemente determinada por las políticas de salario mínimo, la fuerza de las organizaciones sindicales y el grado de formalidad laboral. En el ámbito de los salarios mínimos, investigaciones demuestran \parencite{MaurizioVazquez2016, GindlingRonconi2025} que las políticas de piso salarial en Argentina, Brasil, Chile y Uruguay han ejercido un papel compresor sobre la desigualdad salarial, impactando de forma más acentuada en la parte baja de la distribución. Sin embargo, Firpo y Portella \parencite*{FirpoPortella2025} señalan que en países como Brasil, si bien la brecha de salarios disminuyó significativamente durante ciertos períodos, su persistencia depende estrechamente de la estructura productiva.

La institucionalidad laboral y el poder de negociación colectiva constituyen otro determinante clave. Durán \parencite*{DuranS2026}, Barrera y Noguera \parencite*{BarreraInsuaNoguera2026} analizan los casos de Chile y Argentina respectivamente, mostrando cómo la densidad y el poder sindical afectan directamente la participación de los salarios en el ingreso nacional y reducen la dispersión de las remuneraciones. Por su parte, la precarización y la tipología de contratación imponen restricciones estructurales: Amuedo-Dorantes \parencite*{Amuedo-Dorantes2005} identifica que la heterogeneidad en los contratos de trabajo en Chile genera una marcada brecha de ingresos entre trabajadores permanentes y temporales. Esta precariedad laboral se intensifica en sectores estacionales y exportadores \parencite{TapiaCatalanMarchantSantiago2025}, así como en el ámbito del trabajo irregular o marginado.

Adicionalmente, las dinámicas de inserción laboral femenina e informalidad configuran sesgos estructurales de ingresos. Villanueva y Lin \parencite*{VillanuevaLin2020} demuestran que en América Latina la penalización salarial por maternidad se profundiza sustancialmente bajo esquemas de informalidad laboral. En el caso argentino, Kennedy y Aguila \parencite*{KennedyAguila2021} observan que, si bien la incorporación masiva de las mujeres al mercado de trabajo incrementó los ingresos absolutos de los hogares, la persistencia de la segregación ocupacional contuvo una redistribución más equitativa. En Chile, las investigaciones de la Fundación SOL evidencian que una porción mayoritaria de la fuerza de trabajo percibe salarios insuficientes para cubrir una canasta básica de reproducción \parencite{DuranKremerman2024}, en un contexto atravesado por una severa pobreza de tiempo que afecta de forma desproporcionada a las mujeres \parencite{barrigaMujeresPobrezaTiempo2025}.

\subsection{ii. Dinámicas inflacionarias, shocks macroeconómicos y costo de vida diferenciado}

El impacto del costo de vida sobre la economía doméstica no se distribuye de manera uniforme, pues la inflación afecta con mayor severidad a los estratos de menores ingresos. Adam \parencite*{Adam2025} muestra que en Argentina las tasas de inflación experimentadas por los quintiles más vulnerables superaron de forma sistemática a las de los sectores de altos ingresos. En esta misma línea, do Valle y Gomes da Silva \parencite*{doValleGomesdaSilva2026} confirman para el caso brasileño una alta persistencia y volatilidad en la inflación de la canasta de consumo de los hogares de bajos recursos, lo que debilita progresivamente su capacidad de ahorro. Esta disparidad obliga a los hogares a reconfigurar sus patrones de consumo alimentario \parencite{Nessieretal2026, SilvaCanellaetal2026} hacia bienes de menor valor nutricional o procesados como estrategia de supervivencia.

A nivel macroeconómico, la interacción entre inflación, tipo de cambio y salarios presenta dinámicas complejas según la arquitectura institucional de cada país.

Malher et al. \parencite*{Malheretal2026} y Costa Filho \parencite*{CostaFilho2023} analizan la espiral precio-salario e inflación tendencial en Brasil, destacando el rol del marco de metas de inflación y el traspaso del tipo de cambio. 

En Chile, los determinantes de la inflación y las expectativas del mercado han sido abordados a través del horizonte de política monetaria \parencite{Gredigetal2008}, la persistencia inflacionaria \parencite{Pincheira2008, Medel2025} y el canal de transmisión de shocks económicos hacia la estructura de tasas \parencite{CabellosSanhueza2016, Al-HelalUddinAkram2025}. 

Asimismo, estudios globales sobre economías emergentes sugieren que los componentes globales de inflación imponen límites severos a las políticas de estabilización locales \parencite{Geoetal2024, CakirEkrem2025, Yilmazkuday2025}.

\subsection{iii. Protección social, transferencias monetarias y amortiguación de la pobreza}

Frente a la volatilidad de los ingresos laborales y las presiones inflacionarias, las políticas públicas de transferencias directas constituyen el principal mecanismo estatal de atenuación de la pobreza. Stampini et al. \parencite*{Stampinietal2025} y Brooks \parencite*{Brooks2015} realizan un balance regional del alcance de las transferencias condicionadas y no condicionadas en América Latina, subrayando su efecto positivo inmediato en la reducción de la indigencia, aunque advirtiendo sus limitaciones para alterar la estructura de desigualdad de mercado.

En el plano nacional, las intervenciones presentan trayectorias específicas:
\begin{itemize}
\item \textbf{Chile:} Agostini y Brown \parencite*{AgostiniBrown2011} evalúan el rol de las transferencias monetarias en la reducción de la pobreza regional en Chile, mientras que Oliveria y Osorio \parencite*{OliveiraOsorioGonner2022} analizan las transformaciones e inercias institucionales en la transición desde programas como \textit{Chile Solidario} hacia el \textit{Ingreso Ético Familiar}. Las mediciones del Ministerio de Desarrollo Social y Familia confirman el rol mitigador de estos subsidios en las encuestas CASEN \parencite{ResumenResultadosCASEN2023,TensoresBienestarSocial2024}.
\item \textbf{Argentina:} Carrère \parencite*{Carrere2026} examina la efectividad de la Asignación Universal por Hijo (AUH), demostrando su capacidad para mitigar la inestabilidad de ingresos ante la volatilidad económica y los ciclos inflacionarios.
\item \textbf{Brasil:} La literatura destaca el impacto multifacético de programas como \textit{Bolsa Família}. Segall-Corrêa et al. \parencite*{Segall-Correaetal2008} prueban cómo la transferencia de renta impactó directamente en la disminución de la inseguridad alimentaria en hogares de extrema pobreza, mientras que innovaciones recientes como las monedas sociales y la renta básica local \parencite{OrnellasMauriel2025} han extendido estas redes a escala municipal.
\end{itemize}

Sin embargo, Prieto \parencite*{Prieto2024} y el Informe Final de Recomendaciones para la Actualización de la Medición de la Pobreza elaborada por la Comisión Asesora Presidencial de Expertos y Expertas de Chile \parencite*{InformeFinalRecomendaciones2025} sostienen que la dinámica de bajo ingreso contingente exige una actualización continua de las metodologías de medición y de las estrategias de protección, dado que amplios sectores vulnerables entran y salen constantemente de la línea de pobreza debido a la inestabilidad laboral.

\subsection{iv. Desigualdad estructural: concentración de la riqueza y brechas étnico-raciales}

Los estudios sobre la parte alta de la distribución revelan una alta concentración del ingreso y del capital en los tres países estudiados. Para Chile, los trabajos de Flores et al. \parencite*{Floresetal2020} y Díaz et al. \parencite*{Diazetal2021} demuestran que las encuestas de hogares tradicionales subestiman significativamente los ingresos del 1\% y 0,1\% más rico, requiriendo el uso de registros tributarios y modelos de Pareto ajustados para capturar la real magnitud de la concentración de la riqueza \parencite{Contreras1999, BravoValderrama2011}. En el caso brasileño, Sotomayor \parencite*{Sotomayor2008} y Madeiros y Galvao \parencite*{MadeirosGalvao2014} muestran que la distribución del ingreso ha mostrado una alta inmutabilidad estructural en sus tramos superiores, explicada por el retorno educativo de las élites y la propiedad de activos \parencite{Nascimientoetal1999}.

Por otra parte, la desigualdad socioeconómica latinoamericana se encuentra profundamente racializada y etnificada. Agostini et al. \parencite*{Agostinietal2010} constatan brechas salariales y de pobreza persistentes entre la población indígena y no indígena en Chile. A nivel regional y en el Cono Sur, Painter et al. \parencite*{Painteretal2020}, Sue et al. \parencite*{Sueetal2024}, y Telles et al. \parencite*{Tellesetal2025} demuestran que el tono de piel y la pertenencia étnico-racial constituyen factores determinantes de la desigualdad de activos y riqueza patrimonial, originados históricamente en las estructuras coloniales y esclavistas \parencite{Soaresetal2012, Gradin2014}. En Brasil, Figueiredo Santos \parencite*{FigueiredoSantos2005} y Bailey et al. \parencite*{Baileyetal2014} evidencian que las brechas de ingresos entre la población afrodescendiente y blanca no pueden explicarse únicamente por el nivel educativo, sino que operan mecanismos de discriminación sistémica en el mercado de trabajo.

\subsection{v. Movilidad social, estratificación y percepciones de la desigualdad}

El análisis de la movilidad intergeneracional del ingreso indica una acentuada rigidez en la estructura social de la región. Núñez y Risco \parencite*{NunezRisco2004}, Contreras et al. \parencite*{Contrerasetal2014}, y Daude y Robano \parencite*{DaudeRobano2015} encuentran para el caso chileno y latinoamericano que la posición socioeconómica de los padres predice fuertemente los ingresos de los hijos, limitando el rol de la educación como motor de movilidad ascendente \parencite{Contrerasetal2014}. Resultados análogos son identificados para Argentina \parencite{BeccariaGroisman2008} y Brasil \parencite{CostaRibeiroCarcalhaes2024}.

Esta falta de movilidad y la desconexión entre el crecimiento macroeconómico y el bienestar de los hogares han derivado en tensiones sociopolíticas de gran magnitud. Rojas y Charles-Leija \parencite*{RojasCharlesLeija2022} contrastan la narrativa del "milagro económico" chileno con el estancamiento de los indicadores subjetivos de bienestar. Lattz Lara \parencite*{LattzLara2026} y Pérez-Ahumada et al. \parencite*{Perez-Ahumadaetal2026} analizan cómo la percepción de una desigualdad injusta y la privatización de los servicios básicos constituyeron el detonante del \textit{estallido social} de 2019 en Chile, consolidando una conciencia de clase fuertemente crítica del modelo neoliberal y con niveles heterogéneos de confianza en la acción colectiva y sindical \parencite{Perez-AhumadaCarrasco2025}.